% !TeX program = lualatex
\documentclass[aspectratio=169, 11pt]{beamer}

% ---------------------------------------------------------------------------
% Theme
% ---------------------------------------------------------------------------
\usetheme[
  progressbar=frametitle,
  numbering=fraction,
  block=fill,
  sectionpage=progressbar,
  subsectionpage=none,
]{metropolis}

% ---------------------------------------------------------------------------
% Fonts (requires lualatex or xelatex)
% ---------------------------------------------------------------------------
\usepackage{fontspec}
\setsansfont{Fira Sans}[
  UprightFont = *-Book,
  ItalicFont  = *-BookItalic,
  BoldFont    = *-SemiBold,
]
\setmonofont{Fira Mono}[Scale=MatchLowercase]
\usepackage{unicode-math}
\setmathfont{Fira Math}[Scale=MatchLowercase]

\usepackage[english]{babel}
\usepackage{csquotes}

% ---------------------------------------------------------------------------
% Colors
% ---------------------------------------------------------------------------
\usepackage{xcolor}
\definecolor{accent}{HTML}{2B7A9B}      % teal   - primary accent
\definecolor{accentalt}{HTML}{E8743B}   % orange - secondary / highlights
\definecolor{ink}{HTML}{23373B}         % near-black text
\definecolor{paper}{HTML}{FAFAFA}       % off-white background

\setbeamercolor{normal text}{fg=ink, bg=paper}
\setbeamercolor{alerted text}{fg=accentalt}
\setbeamercolor{example text}{fg=accent}
\setbeamercolor{progress bar}{fg=accent}
\setbeamercolor{title separator}{fg=accent}
\setbeamercolor{progress bar in head/foot}{fg=accent}
\setbeamercolor{progress bar in section page}{fg=accent}
\setbeamercolor{frametitle}{bg=ink, fg=paper}

% ---------------------------------------------------------------------------
% Packages
% ---------------------------------------------------------------------------
\usepackage{booktabs}
\usepackage{graphicx}
\usepackage{amssymb}
\usepackage{pifont}
\newcommand{\cmark}{\ding{51}}
\newcommand{\xmark}{\ding{55}}
\graphicspath{{figures/}}
\usepackage{tikz}
\usetikzlibrary{arrows.meta, positioning, calc, fit, backgrounds, decorations.pathreplacing}
\usepackage{pgfplots}
\pgfplotsset{compat=1.18}
\usepackage{appendixnumberbeamer}
\usepackage{fontawesome5}

% ---------------------------------------------------------------------------
% Code listings
% ---------------------------------------------------------------------------
\usepackage{listings}
\definecolor{codebg}{HTML}{F2F2F0}
\definecolor{codekw}{HTML}{2B7A9B}
\definecolor{codestr}{HTML}{6A8E3B}
\definecolor{codecom}{HTML}{9A9A9A}

\lstdefinestyle{talk}{
  backgroundcolor=\color{codebg},
  basicstyle=\ttfamily\small,
  keywordstyle=\color{codekw}\bfseries,
  stringstyle=\color{codestr},
  commentstyle=\color{codecom}\itshape,
  numberstyle=\ttfamily\scriptsize\color{codecom},
  showstringspaces=false,
  breaklines=true,
  breakatwhitespace=true,
  columns=fullflexible,
  keepspaces=true,
  frame=none,
  xleftmargin=6pt,
  xrightmargin=6pt,
  framexleftmargin=6pt,
  framexrightmargin=6pt,
  aboveskip=1em,
  belowskip=1em,
  upquote=true,
  escapeinside={(*@}{@*)},
}
\lstset{style=talk, language=Python}

\lstnewenvironment{pycode}[1][]{\lstset{language=Python,#1}}{}
\lstnewenvironment{shcode}[1][]{\lstset{language=bash,#1}}{}
\newcommand{\py}[1]{\lstinline[language=Python]|#1|}

% ---------------------------------------------------------------------------
% Small conveniences
% ---------------------------------------------------------------------------
\newcommand{\hl}[1]{\textcolor{accentalt}{\textbf{#1}}}

% Tighter, nicer lists
\setbeamertemplate{itemize items}[circle]
\setbeamertemplate{itemize subitem}[triangle]

% Frame with a full-bleed image background:
%   \fullimageframe{myfigure}{Optional caption}
\newcommand{\fullimageframe}[2]{%
  \begin{frame}[plain]
    \begin{tikzpicture}[remember picture, overlay]
      \node[at=(current page.center)] {%
        \includegraphics[width=\paperwidth, height=\paperheight, keepaspectratio]{#1}};
    \end{tikzpicture}
    \ifx\relax#2\relax\else\vfill{\small #2}\fi
  \end{frame}%
}


\title{MLSS 2026 Tutorial}
\subtitle{Efficient LLM Inference and Serving}
\author{Alexander Conzelmann and Shiwei Liu}
\date{\today}
% \institute{}

\begin{document}

\maketitle

\begin{frame}{Outline}
  \textbf{Goal:} Get practical intuition on how GPUs work and how to use vLLM. 
  \begin{itemize}
    \item Prefill vs Decode: GPU memory efficiency 
    \item Speculative Decoding: Speeding up inference losslessly
  \end{itemize}
\end{frame}

\begin{frame}[standout]
  Prefill versus Decode \\ \small How GPUs handle memory and computation
\end{frame}

\begin{frame}{Two Distinct Stages in LLM Inference}
Due to KV-caching, inference can be separated into \textbf{Prefill} and \textbf{Decoding}. Let us assume
we are working on a single prompt at a time:
%These stages have

\textbf{Prefill.} The tokenized user prompt $P = s_1 s_2 \ldots s_m$ is processed. Output is the next token 
$\hat{s}_{m+1}$ plus per-layer $l$ and per-token $k$ KV cache values $K_{l,k}, V_{l,k}$.

Operations on hidden states $X$, such as  calculating $Q,K,V$ values are $(m, d) \times (d,d)$.

\textbf{Decode.} We calculate one output token using the previous prompts' cached KV values. 

Operations on hidden states $x$, such as  calculating $Q,K,V$ values are $({\color{red}1}, d) \times (d,d)$.
\end{frame}

% Insert after "Two Distinct Stages in LLM Inference"
\begin{frame}{What Actually Gets Computed}
Per layer, on hidden states $X$, with $W_Q \in \mathbb{R}^{d \times H_q h}$, $W_K, W_V \in \mathbb{R}^{d \times H_{kv} h}$:
\[
  Q = X W_Q, \quad K = X W_K, \quad V = X W_V,
  \qquad
  \mathrm{Attn}(Q,K,V) = \mathrm{softmax}\!\left(\frac{QK^\top}{\sqrt{h}}\right) V
\]
In prefill $X \in \mathbb{R}^{m \times d}$. In decode $X = x \in \mathbb{R}^{{\color{red}1} \times d}$, and $K, V$ are the
cached ones with the single new row appended.

\end{frame}

\begin{frame}{What Actually Gets Computed}
\vspace{0.3em}
\textbf{Qwen3-4B:} $d = 2560$, $L = 36$ layers, $H_q = 32$ query heads, $H_{kv} = 8$ KV heads,
head dim $h = 128$, $d_{ff} = 9728$.

\vspace{0.3em}
\begin{table}
\centering
\footnotesize
\begin{tabular}{lcc}
\toprule
per layer & \textbf{Prefill} ($m$ tokens) & \textbf{Decode} (1 token, cache $m$) \\
\midrule
$Q$ proj            & $(m,d)\times(d,H_q h)$    & $({\color{red}1},d)\times(d,H_q h)$ \\
$K,V$ proj          & $(m,d)\times(d,H_{kv}h)$  & $({\color{red}1},d)\times(d,H_{kv}h)$ \\
scores $QK^\top$    & $(m,h)\times(h,m)$        & $({\color{red}1},h)\times(h,m)$ \\
$\mathrm{softmax}\cdot V$ & $(m,m)\times(m,h)$  & $({\color{red}1},m)\times(m,h)$ \\
MLP                 & $(m,d)\times(d,d_{ff})$   & $({\color{red}1},d)\times(d,d_{ff})$ \\
\midrule
FLOPs               & $\mathcal{O}(m d^2 + m^2 d)$ & $\mathcal{O}(d^2 + m d)$ \\
\bottomrule
\end{tabular}
\end{table}
\textbf{Problem:} The cache removes the compute, but not the weight reads. Every matrix on the right
still has to be pulled from HBM in full, to be multiplied against a single row.
\end{frame}

% ---------------------------------------------------------------------------
% GPU memory hierarchy (A100 SXM 80GB) -- motivates the roofline frame
% ---------------------------------------------------------------------------
\begin{frame}{The GPU Memory Hierarchy}
  \vspace{-0.4em}
  \begin{center}
  \resizebox{\textwidth}{!}{%
  \begin{tikzpicture}[
      font=\scriptsize,
      lvl/.style={
        draw=ink!35, line width=0.5pt, rounded corners=2.5pt,
        minimum height=0.82cm, text=ink, align=center, inner sep=2pt},
      note/.style={text width=3.5cm, align=left, font=\scriptsize, text=ink!75},
      brc/.style={decorate, decoration={brace, amplitude=4pt}, draw=ink!45},
    ]

    % ---- the levels -------------------------------------------------------
    \node[lvl, minimum width=4.2cm, fill=accentalt!85, text=white] (reg)  at (0, 0.0)
      {\textbf{Registers} \;\textbar\; 256\,KB/SM \;\textbar\; $\sim$100\,TB/s};
    \node[lvl, minimum width=5.6cm, fill=accentalt!45] (smem) at (0,-1.0)
      {\textbf{L1 / Shared Memory} \;\textbar\; 192\,KB/SM \;\textbar\; $\sim$19\,TB/s};
    \node[lvl, minimum width=7.0cm, fill=accent!35] (l2)   at (0,-2.0)
      {\textbf{L2 Cache} \;\textbar\; 40\,MB (chip-wide) \;\textbar\; $\sim$5\,TB/s};
    \node[lvl, minimum width=8.4cm, fill=accent!70, text=white] (hbm)  at (0,-3.0)
      {\textbf{HBM2e -- global memory} \;\textbar\; 80\,GB \;\textbar\; \textbf{2039\,GB/s}};
    \node[lvl, minimum width=9.8cm, fill=ink!8, draw=ink!25, dashed] (off) at (0,-4.15)
      {\textbf{Off-chip} \;\textbar\; NVLink 600\,GB/s \;\textbar\; PCIe\,4.0 64\,GB/s \;\textbar\; host DRAM};

    % ---- left gradient arrow ---------------------------------------------
    \draw[-{Stealth[length=4pt]}, ink!45, line width=0.8pt]
      (-5.35, 0.45) -- (-5.35,-4.6);
    \node[rotate=90, anchor=south, font=\tiny, text=ink!60] at (-5.5,-2.0)
      {capacity $\uparrow$ \quad latency $\uparrow$ \quad bandwidth $\downarrow$};

    % ---- right annotations ------------------------------------------------
    \draw[brc] (5.25, 0.42) -- (5.25,-1.42);
    \node[note, anchor=west] at (5.6,-0.5)
      {\textbf{Where reuse happens.} Data is pulled here and multiplied many times.
       Arithmetic intensity counts exactly this: FLOPs per byte lifted out of HBM.};

    \draw[brc] (5.25,-2.62) -- (5.25,-3.42);
    \node[note, anchor=west] at (5.6,-3.02)
      {\textbf{Where the model lives.} 8\,GB of weights $+$ the KV cache.
       Every decode step re-reads all of it.};

  \end{tikzpicture}}
  \end{center}

  \vspace{-0.2em}
  Only the bottom level is big enough to hold the model, but it is very slow. 
  For matrix multiplications, the weight matrices need to be pulled into higher GPU memory levels.
\end{frame}


\begin{frame}{Roofline Model: Compute- vs Memory-Bound}
A GPU operation is bound by whichever resource runs out first: \textbf{compute} (FLOPs/s)
or \textbf{memory bandwidth} (bytes/s). The relevant quantity is \textbf{arithmetic intensity}:
\[
  \text{AI} = \frac{\text{FLOPs}}{\text{Bytes moved from HBM}} \quad \left[\frac{\text{FLOP}}{\text{byte}}\right]
\]
Achieved throughput is capped by the \textbf{roofline}:
\[
  \text{Achieved FLOP/s} = \min\big(\text{Peak Compute},\ \text{AI} \times \text{Memory Bandwidth}\big)
\]
The \textbf{ridge point} is where the two terms are equal --- below it, you are memory-bound
no matter how fast your cores are.

\vspace{0.5em}
\textbf{Reference GPU (A100 SXM 80GB, FP16 tensor cores):}
\[
  \text{Peak Compute} = 312~\text{TFLOP/s}, \quad
  \text{Bandwidth} = 2039~\text{GB/s}
  \ \Rightarrow\ \text{Ridge} \approx 153~\tfrac{\text{FLOP}}{\text{byte}}
\]
Any matmul with AI below $153$ leaves compute cores idle, waiting on HBM.
\end{frame}

\begin{frame}{Why Decoding is Slow}
Take one linear projection of \textbf{Qwen3-4B}: $d = 2560$ 

\vspace{0.3em}
\begin{table}
\centering
\small
\begin{tabular}{lcc}
\hline
 & \textbf{Decode} $(1,d)\times(d,d)$ & \textbf{Prefill} $(m{=}512,d)\times(d,d)$ \\
\hline
FLOPs & $2d^2 \approx 13.1$ MFLOP & $2md^2 \approx 6.7$ GFLOP \\
Bytes moved (weights) & $2d^2 \approx 13.1$ MB & $2d^2 \approx 13.1$ MB \\
Arithmetic Intensity & $\mathbf{1}$ FLOP/byte & $\mathbf{512}$ FLOP/byte \\
Achieved throughput & $2.0$ TFLOP/s & $312$ TFLOP/s \\
GPU utilization & $\mathbf{0.65\%}$ & $\mathbf{100\%}$ \\
\hline
\end{tabular}
\end{table}

$\implies$ during decoding, the GPUs sit almost completely idle.

\textbf{Batching} can help, but only if you have multiple requests 
in parallel you can batch together.
\end{frame}


% Insert after "Why Decoding is Slow"
% \begin{frame}{Where the Bytes Come From}
% A decode step reads two things: \textbf{all weights} and \textbf{the entire KV cache}.
%  
% \textbf{Weights.} $4$B params in FP16 $=8$ GB, independent of $m$.
%  
% \textbf{KV cache.} Per token: $2 L H_{kv} h \cdot 2\,\mathrm{B} = 2\cdot36\cdot8\cdot128\cdot2 \approx 147$ KB.
% At $m = 8192$: $\approx 1.2$ GB --- and it grows every step.
%  
% \vspace{0.3em}
% \begin{table}
% \centering
% \small
% \begin{tabular}{lccc}
% \toprule
% & \textbf{FLOPs} & \textbf{Bytes} & \textbf{AI} \\
% \midrule
% weight matmuls ($N$ params) & $2N$            & $2N$              & $1$ \\
% attention over cache        & $4 m L H_q h$   & $4 m L H_{kv} h$  & $H_q / H_{kv} = 4$ \\
% \bottomrule
% \end{tabular}
% \end{table}
%  
% Both sit far below the ridge point of $153$; note the attention AI is $m$-independent, GQA is the only
% thing lifting it above $1$.
% \[
%   T_{\text{step}} \geq \frac{(8 + 1.2)~\text{GB}}{2039~\text{GB/s}} \approx 4.5~\text{ms}
%   \quad\Rightarrow\quad \approx 220~\text{tok/s ceiling}
% \]
%  
% \textbf{Batching only fixes half of it.} $B$ requests share the weights (AI $\to B$), but each carries its
% own cache: attention AI stays at $4$ and cache traffic scales with $B$.
% \end{frame}

\begin{frame}[standout]
  DEMO
\end{frame}

\begin{frame}[standout]
  Speculative Decoding \\ \small How to speed up our model (almost) for free
\end{frame}

\begin{frame}{Speculative Decoding}
  % Sometimes, you cannot batch requests together (i.e. local, offline LLMs). One way to improve arithmetic intensity during decoding
  \textbf{Idea.} A cheap model (drafter) predicts tokens, original LLM (verifier) confirms tokens. 

  During decoding, drafter produces $b$ proposal tokens $\hat{s}_{m+1}, \hat{s}_{m+2}, \ldots, \hat{s}_{m+b}$. 
  Verifier then batches the proposal tokens together to determine logits at each proposal position.

  Where top-1 logit agrees with proposal: accept token, else reject. After the first rejected position, 
  the full draft is thrown away 
  $\implies$ the verifiers probability distribution is unaltered (lossless speedup).

\end{frame}

\begin{frame}{Why does speculative decoding speed up inference?}
  Roofline argument: in memory-bound regime, processing $M$ tokens is roughly as expensive as processing $1$ token
  {\color{gray}(until cross-over: so $M_{\text{cross}} \approx 153$ on the A100)}.
  % Sometimes, you cannot batch requests together (i.e. local, offline LLMs). One way to improve arithmetic intensity during decoding

\textbf{Per-round timing.}
\[
  T_{\text{round}} = \gamma T_d + T_v
\]
where $T_v$ = verifier forward-pass time, $T_d$ = drafter forward-pass time

\textbf{Expected tokens per round.} With per-position acceptance rate $\alpha_j$: 
$$E[\text{tokens/round}] = 1 + \sum_{i=1}^\gamma \prod_{j=1}^i \alpha_j \stackrel{\text{if} \alpha_j = \alpha \forall j}{=} 1 + \sum_{i=1}^\gamma \alpha^i = \frac{1-\alpha^{\gamma+1}}{1-\alpha}$$

\textbf{Speedup} vs.\ plain autoregressive decoding (1 token every $T_v$):
\[
  \boxed{\text{Speedup} = \frac{T_v}{\gamma T_d + T_v}\cdot\frac{1-\alpha^{\gamma+1}}{1-\alpha}}
\]
\end{frame}

\begin{frame}{N-Gram Speculation}
  \textbf{Idea.} Take the last $n$ tokens, search for the same $n$-gram
  earlier in prompt $+$ generation, propose the $\gamma$ tokens that followed it there.
  Verification is unchanged $\implies$ still lossless.

  \textbf{Cost.} String lookup is almost free: $T_d \approx 0$.
  \[
    \text{Speedup} = \frac{T_v}{\gamma T_d + T_v}\cdot\frac{1-\alpha^{\gamma+1}}{1-\alpha}
    \quad\xrightarrow{\ T_d \to 0\ }\quad
    \frac{1-\alpha^{\gamma+1}}{1-\alpha}
  \]
  $\implies$ no interior optimum in $\gamma$ any more; only the verifier batch approaching
  $M_{\text{cross}}$ still limits it.

  \textbf{Where it helps.} $\alpha$ is purely lexical: a draft is right only if the continuation
  \emph{literally} stands in the context. Chat: $\alpha \approx 0$.
  Summarization, RAG, code edits: high.

  \textbf{Free:} no second model, no weights in HBM, no shared tokenizer, no training.
\end{frame}

\begin{frame}{Other Drafters}
  In principal, we can use any method that calculates new possible tokens given previous text 
  and possibly verifier information like hidden states. Popular methods:

  \begin{table}
  \centering
  \small
  \begin{tabular}{lccc}
  \toprule
  \textbf{Drafter} & $T_d$ & $\alpha$ & \textbf{needs training} \\
  \midrule
  standalone model    & $0.1$--$0.2\,T_v$        & moderate       & \xmark \\
  EAGLE / Medusa head & $\ll T_v$                & high           & \cmark \\
  Diffusion LM        & one pass for all $\gamma$ & moderate      & \cmark \\
  $n$-gram lookup     & $\approx 0$              & task-dependent & \xmark \\
  \bottomrule
  \end{tabular}
  \end{table}

  \textbf{EAGLE.} Drafts from the verifier's hidden states instead of from text $\implies$ high $\alpha$
  for one extra layer. Proposes a \emph{tree} of continuations, verified in one pass.

  \textbf{Diffusion.} All $\gamma$ tokens in a single bidirectional pass: $\gamma T_d \to T_d$, so the
  term that created the interior optimum in $\gamma$ disappears.
\end{frame}

\begin{frame}[standout]
  DEMO
\end{frame}

\begin{frame}[standout]
  Quantization \\ \small How to get even faster, at a cost
\end{frame}

\begin{frame}{Quantization}
  Quantization is a \textit{lossy} speed-up technique: we perform matrix multiplications in lower precision. 

  \textbf{Quantization does not always increase throughput.} Some formats are supported by dedicated 
  tensor cores on GPU, only these are performed faster. 
  Else: very nuanced, faster memory transfer but slower execution (due to dequantization overhead).
  \begin{table}[htbp]
\centering
\caption{Quantization kernel support by NVIDIA GPU architecture}
\label{tab:quant-support}
\begin{tabular}{@{}lccc@{}}
\toprule
\textbf{Architecture} & \textbf{INT8} & \textbf{FP8} & \textbf{FP4} \\
\midrule
Ampere    & \cmark & \xmark & \xmark \\
Hopper    & \cmark & \cmark & \xmark \\
Blackwell & (\cmark) & \cmark & \cmark \\
\bottomrule
\end{tabular}
\end{table}
\end{frame}



\begin{frame}{RTN: Uniform Affine Quantization}
  We approximate a float vector by \textbf{integers times a float scale}: $\hat{w} = s\,(w^{\text{int}} - z)$.

  \textbf{Asymmetric} (scale $s$, zero point $z$), grid spans $[\min w, \max w]$:
  \[
    s = \frac{\max(w) - \min(w)}{2^b - 1},
    \quad
    z = \left\lfloor -\frac{\min(w)}{s} \right\rceil,
    \quad
    w^{\text{int}} = \mathrm{clamp}\!\left(\left\lfloor \frac{w}{s} \right\rceil + z,\ 0,\ 2^b-1\right)
  \]

  \textbf{Symmetric} drops $z$, grid is centred on zero:
  \[
    s = \frac{\max |w|}{2^{b-1}-1},
    \qquad
    w^{\text{int}} = \mathrm{clamp}\!\left(\left\lfloor \frac{w}{s} \right\rceil,\ -2^{b-1},\ 2^{b-1}-1\right)
  \]

  \textbf{Round-to-nearest} = do exactly this, no data, no search. One pass over the checkpoint.

  Two error sources, traded against each other by $s$: \hl{rounding} error $\leq s/2$ per weight, and
  \hl{clipping} error for everything outside the grid. A single outlier inflates $s$ and thereby
  degrades \emph{every} weight sharing it.
\end{frame}

\begin{frame}{Granularity: What Shares a Scale}
  The only real knob in RTN is how large a group shares one $(s, z)$.

  \vspace{-0.2em}
  \begin{table}
  \centering
  \footnotesize
  \begin{tabular}{lccc}
  \toprule
  \textbf{granularity} & \textbf{scales/matrix} & \textbf{extra bits/weight} & \textbf{robustness} \\
  \midrule
  per-tensor            & $1$                   & $\approx 0$        & worst \\
  per-channel (per row) & $d_{\text{out}}$      & $\approx 0$        & better \\
  group-$g$ ($g{=}128$) & $d_{\text{out}} d_{\text{in}} / g$ & $32/g = 0.25$ & best \\
  \bottomrule
  \end{tabular}
  \end{table}
  \vspace{-0.2em}

  \textbf{Effective bits} $= b + 32/g$ (fp16 scale $+$ zero point per group): 4-bit group-128
  really costs $4.25$ bit/weight.

  \textbf{Empirically.} 8 bit is free at any granularity. 4 bit \emph{needs} groups: per-tensor is
  broken, group-128 is within a few percent. 3 bit is the edge, 2 bit is gone.

  $\implies$ every production 4-bit format (AWQ, GPTQ, NF4) is group-wise.
\end{frame}

\begin{frame}{Dequantization: Where the Integers Turn Back Into Floats}
  Storing $w^{\text{int}}$ is only half the story --- the matmul has to consume it somehow.

  \textbf{W4A16 (AWQ, GPTQ).} Weights 4-bit in HBM, activations fp16. No 4-bit tensor core exists,
  so the kernel \hl{dequantizes in registers} and multiplies in fp16:
  \[
    y = x \hat{W}^{\mathsf{T}}, \qquad \hat{W} = s \odot (W^{\text{int}} - z) \ \text{unpacked on the fly}
  \]
  Win: $4\times$ fewer weight bytes. Cost: unpack $+$ scale per weight, every matmul, every step.

  \textbf{W8A8 (INT8, FP8).} Both operands low precision $\implies$ the tensor core does the matmul
  natively; dequantization is one multiply on the \emph{output}, not on the weights.

  \textbf{Roofline consequence.} W4A16 lifts decode AI from $1$ to $\approx 4$ FLOP/byte and adds
  dequant work. Batch 1: near-$4\times$. Large batch: frequently \emph{slower} than bf16.
\end{frame}

\begin{frame}{Accumulation}
  A MAC array loads accumulators with the bias and adds products of the loaded operands:
  \[
    \hat{A}_n = \hat{b}_n + \sum_m \hat{W}_{n,m}\, \hat{x}_m
    = \hat{b}_n + \sum_m \left(s_w W^{\text{int}}_{n,m}\right)\!\left(s_x x^{\text{int}}_m\right)
    = \hat{b}_n + \hl{$s_w s_x$} \sum_m W^{\text{int}}_{n,m} x^{\text{int}}_m
  \]
  The scales \textbf{factor out of the sum} $\implies$ the inner loop is pure integer arithmetic.

  \textbf{With zero points} the cross terms do not vanish:
  \[
    \sum_m (W^{\text{int}}_{n,m} - z_w)(x^{\text{int}}_m - z_x)
    = \sum_m W^{\text{int}}_{n,m} x^{\text{int}}_m
    - z_w \!\underbrace{\sum_m x^{\text{int}}_m}_{\text{runtime}}
    - z_x \!\underbrace{\sum_m W^{\text{int}}_{n,m}}_{\text{offline}}
    + \underbrace{d_{\text{in}} z_w z_x}_{\text{offline}}
  \]
  $\implies$ symmetric weights ($z_w = 0$) kill the one term that costs anything at runtime. This is why
  weight quantization in practice is symmetric even though it wastes a level.
\end{frame}

\begin{frame}{Accumulator Width and Requantization}
  \textbf{Products do not fit.} INT8 $\times$ INT8 needs $16$ bit; summing $d_{\text{in}}$ of them needs
  $16 + \lceil \log_2 d_{\text{in}} \rceil$ bit --- for $d = 2560$ that is $28$ bit $\Rightarrow$ INT32.
  Overflow is silent, so accumulators are always wider: INT8$\to$INT32, FP8$\to$FP32, fp16$\to$fp32.

  \textbf{Requantization.} The INT32 accumulator must become an INT8 activation before the next layer,
  or you gained nothing on bandwidth:
  \[
    x^{\text{int}}_{\text{next}} = \mathrm{clamp}\!\left(\left\lfloor \frac{s_w s_x}{s_y} A^{\text{int}}_n \right\rceil + z_y,\ 0,\ 255\right)
  \]
  with $s_y$ from activation calibration. Fusing this into the matmul epilogue is what makes W8A8 fast.

  \textbf{Why W8A8 is not group-wise.} Group scales break the factorization: $s_w$ varies \emph{inside}
  the sum, so the accumulator must be split per group and rescaled. Group-128 therefore lives in
  W4A16, where you dequantize anyway.
\end{frame}

\begin{frame}{GPTQ: Rounding Is the Wrong Objective}
  RTN minimizes \emph{weight} error $\|W - \hat{W}\|$. What we actually care about is \emph{layer output}
  error on real data. With calibration activations $X \in \mathbb{R}^{d_{\text{in}} \times N}$:
  \[
    \boxed{\ \hat{W}^\ast = \arg\min_{\hat{W}} \ \big\| W X - \hat{W} X \big\|_2^2 \ }
  \]
  Separable across output \emph{rows} $\implies$ solve each row independently. \emph{Not} separable
  across input columns: a direction $X$ excites often deserves a smaller error. That coupling is
  exactly what RTN throws away.

  \textbf{Key quantity.} The Hessian w.r.t.\ one row is the input covariance
  $H = 2 X X^{\mathsf{T}} \in \mathbb{R}^{d_{\text{in}} \times d_{\text{in}}}$ --- one forward pass over
  $\approx 128$ calibration sequences, accumulated per layer.
\end{frame}

\begin{frame}{GPTQ: The Update}
  Quantize columns left to right. Having fixed column $q$ to its grid value, the \textbf{optimal
  correction} to the remaining (still-unquantized) columns $F$ is the OBS update:
  \[
    e_q = \frac{w_q - \mathrm{quant}(w_q)}{[H_F^{-1}]_{qq}},
    \qquad
    \delta_F = -\,e_q \cdot \big(H_F^{-1}\big)_{q,\,F}
  \]
  $\implies$ every rounding error is \hl{immediately compensated} by nudging weights we have not
  committed to yet. Errors get absorbed instead of accumulating.

  \textbf{Why it is tractable.} OBQ uses a per-row greedy order and re-inverts $H$ each step:
  $\mathcal{O}(d_{\text{in}}^4)$ per row. GPTQ fixes \emph{one} column order for all rows $\implies$ a
  single Cholesky of $H^{-1}$, all rows updated as one matrix op.

  \textbf{Practicalities.} Damping $H \leftarrow H + \lambda \bar{d} I$ ($\lambda \approx 1\%$), else the
  Cholesky fails. \emph{Act-order}: columns by descending $H_{qq}$, so sensitive ones are quantized
  while correction budget remains. Lazy block updates ($128$ columns) to stay compute-bound.
\end{frame}

\begin{frame}{Where This Lands}
  \begin{table}
  \centering
  \small
  \begin{tabular}{lcccc}
  \toprule
  & \textbf{data} & \textbf{cost} & \textbf{3--4 bit} & \textbf{8 bit} \\
  \midrule
  RTN   & none                & seconds                     & breaks       & fine \\
  GPTQ  & $\approx 128$ seqs  & minutes--hours, $\mathcal{O}(d^3)$/layer & good & fine \\
  AWQ   & $\approx 128$ seqs  & minutes                     & good         & fine \\
  \bottomrule
  \end{tabular}
  \end{table}

  \textbf{AWQ, in one line.} Instead of correcting rounding errors, \emph{rescale} channels before
  rounding: $\hat{W} = Q(W \cdot \mathrm{diag}(\alpha))$, $\hat{x} = \mathrm{diag}(\alpha)^{-1} x$,
  with $\alpha$ chosen from activation magnitudes so salient channels land on a finer effective grid.
  Mathematically free at inference --- $\alpha$ folds into the previous layer.

  \vspace{0.3em}
  $\implies$ all three are \emph{indistinguishable} at 8 bit. They only separate in the 3--4 bit range,
  which is precisely where a $2\times$ over FP8 is on the table.
\end{frame}

% ---------------------------------------------------------------------------
% Content sections go here, e.g.:
%   \section{Introduction}
%   \input{sections/01-introduction}
% ---------------------------------------------------------------------------

\appendix

\begin{frame}[standout]
  Questions?
\end{frame}

\end{document}
